\section*{Aufgabenstellung}
Gibt es Spuren im Serverimage bzgl.
der Installation des Tor Dienstes?
Gibt es einen Webauftritt?

\section*{Lösung}
Da das Image verschlüsselt ist (siehe Aufgabe 2), erfolgt die Analyse auf dem raw image.
Hierfür wurde das Image zuerst mit ewfmount gemountet:

\begin{lstlisting}[language=bash]
mkdir -p /mnt/ewf_server
ewfmount images/20260328_server.E01 /mnt/ewf_server
\end{lstlisting}

Mithilfe von strings konnte ich die folgenden Onion Domain und Fragmente des Apache Webservers finden:
\begin{lstlisting}[language=bash]
root@Alpha-17:/home/luis/Programming/GitHub/University/DigitalForensic# strings /mnt/ewf_server/ewf1 | grep -i '\.onion' | sort -u
Homepage: http://2xzrgzkwdqvcfw2ckas4n6iu4gcyk3mvfh6eq5zt5il3ohz2qbtvssyd.onion/
Not resolving .onion address (RFC 7686)
.onion
.onion.
root@Alpha-17:/home/luis/Programming/GitHub/University/DigitalForensic# strings /mnt/ewf_server/ewf1 | grep -i 'apache\|nginx\|DocumentRoot' | sort -u | head -10
                    $APACHE2CTL $METH > /dev/null 2>&1
                $APACHE2CTL graceful > /dev/null 2>&1
                $APACHE2CTL start
        [ "$APACHE2_MAINTSCRIPT_NAME" ] || APACHE2_MAINTSCRIPT_NAME="${0##*.}"
                                    [ "$APACHE2_MAINTSCRIPT_NAME" = 'postinst' ] || [ "$APACHE2_MAINTSCRIPT_NAME" = 'preinst' ] ; then
  . $APACHE_ENVVARS
      . $apache_reload
    $apache_reload = "  service apache2$dir_suffix $reload\n";
    $apache_reload = "  systemctl $reload $apache_service\n";
    $apache_service .= $service_suffix
\end{lstlisting}

Die Installation von Tor konnte jedoch nicht direkt nachgewiesen werden:
\begin{lstlisting}[language=bash]
root@Alpha-17:/home/luis/Programming/GitHub/University/DigitalForensic# strings /mnt/ewf_server/ewf1 | grep -i 'apt.*install.*tor' | head -10
 _ZN3APT11CacheFilter19PackageIsNewInstallclERKN8pkgCache11PkgIteratorE@APTPKG_7.0 1.1~exp4
 _ZN3APT12StateChanges7InstallERKN8pkgCache11VerIteratorE@APTPKG_7.0 1.1~exp15
 _ZN3APT14CacheSetHelper22canNotFindInstalledVerER12pkgCacheFileRKN8pkgCache11PkgIteratorE@APTPKG_7.0 0.8.0
 _ZN3APT25VersionContainerInterface15getInstalledVerER12pkgCacheFileRKN8pkgCache11PkgIteratorERNS_14CacheSetHelperE@APTPKG_7.0 0.8.16~exp9
_ZN3APT11CacheFilter19PackageIsNewInstallclERKN8pkgCache11PkgIteratorE
_ZN3APT14CacheSetHelper22canNotFindInstalledVerER12pkgCacheFileRKN8pkgCache11PkgIteratorE
_ZN3APT25VersionContainerInterface15getInstalledVerER12pkgCacheFileRKN8pkgCache11PkgIteratorERNS_14CacheSetHelperE
_ZN3APT12StateChanges7InstallERKN8pkgCache11VerIteratorE
virtual bool APT::Internal::Patterns::PackageIsInstalled::operator()(const pkgCache::PkgIterator&)
virtual bool APT::Internal::Patterns::PackageIsInstalled::operator()(const pkgCache::VerIterator&)
\end{lstlisting}

Nach der analyse des squashfs Containers aus Aufgabe 8 würde ich sagen der Linux Server definitiv für Darknet Aktivitäten genutzt wurde.
Ein Web Server war aber vermutlich nicht installiert, da es keine Spuren von Apache oder Nginx gibt.

\begin{lstlisting}[language=bash]
root@Alpha-17:/home/luis/Programming/GitHub/University/DigitalForensic# strings /tmp/swap.bin | grep -i '\.onion' | sort -u
root@Alpha-17:/home/luis/Programming/GitHub/University/DigitalForensic# strings /tmp/swap.bin | grep -i 'tor\|hidden_service' | sort -u | head -20
root@Alpha-17:/home/luis/Programming/GitHub/University/DigitalForensic# strings /tmp/swap.bin | grep -i 'apache\|nginx\|www' | sort -u | head -20
root@Alpha-17:/home/luis/Programming/GitHub/University/DigitalForensic# binwalk /mnt/ewf_server/ewf1 2>&1 | grep -i '.onion|tor|hidden_service|apache|nginx|www'
root@Alpha-17:/home/luis/Programming/GitHub/University/DigitalForensic#
\end{lstlisting}



